\subsection{A finite determinantal test for special hyperplanes}

\begin{proposition}[A2D; Proved]\label{a:determinantal}
Let $X/k$ be smooth projective and integral of dimension at least two,
where $k$ is algebraically closed of characteristic $p>0$.
Let $L$ be very ample, let $M\ne\OO_X$ satisfy $M^p\simeq\OO_X$,
and put
\[
 a=h^1(X,M\otimes L^{-1}),\qquad b=h^1(X,M),\qquad
 r=h^0(X,L)-1.
\]
On the open subset of $|L|=\mathbb P^r$ parametrizing smooth divisors,
the locus where $M|_Y$ is trivial is finite. It is the restriction of
the maximal-minor locus of the linear multiplication matrix
\[
 H^1(X,M\otimes L^{-1})\otimes\OO_{\mathbb P^r}(-1)
 \longrightarrow H^1(X,M)\otimes\OO_{\mathbb P^r}.
\]
We have $a\leq b$. If this finite locus is nonempty, then $a>0$ and
\begin{equation}\label{a:determinantal-bound}
 b-a\ \geq\ h^0(X,L)-2.
\end{equation}
In dimension at least three, it follows that restriction on geometric
Picard groups is injective for every smooth member of $|L|$ if, for
each nontrivial $M\in\Pic(X)[p]$, either $a=0$ or the reverse strict
inequality to \eqref{a:determinantal-bound} holds.
\end{proposition}

\begin{proof}
The injective logarithmic differential construction of
Proposition~\ref{a:generic} gives a nonzero form
$\eta=d\log(M)\in H^0(X,\Omega_X^1)$.
If $M|_Y$ is trivial, its restriction to $\Omega_Y^1$ is zero.
The conormal sequence and
$H^0(Y,\OO_Y(-Y))=0$ therefore show that
$\eta|_Y=0$ already in $\Omega_X^1|_Y$.
Every such $Y$ is consequently a divisorial component of the zero
scheme of $\eta$. A nonzero section of a vector bundle on an integral
noetherian scheme has only finitely many divisorial components in its
zero scheme. A smooth ample divisor here is connected and hence
integral. There are therefore only finitely many possible $Y$.

A nontrivial numerically trivial line bundle on a projective integral
scheme has no nonzero section: any nonzero effective divisor has
positive intersection with a suitable power of an ample divisor.
The same intersection argument gives $H^0(X,M\otimes L^{-1})=0$.
For a section $h$ cutting out a smooth $Y$, the divisor sequence gives
\[
 H^0(Y,M|_Y)=\ker\bigl(H^1(X,M\otimes L^{-1})
                  \xrightarrow{\,h\,}H^1(X,M)\bigr).
\]
Since $M|_Y$ is numerically trivial, its space of sections has dimension
one exactly when it is trivial, and dimension zero otherwise.
This proves the asserted determinantal description. There are smooth
members outside the finite exceptional locus, so a multiplication map
is injective and $a\leq b$.

At a bad smooth member, the matrix has rank exactly $a-1$.
If $a=0$ there is no such member. Otherwise choose an invertible
$(a-1)$-minor locally, allowing the empty minor when $a=1$.
Elementary row and column operations reduce the maximal-minor ideal
locally to $b-a+1$ functions, the entries of its remaining column.
The bad locus has dimension zero in the smooth-member open subset
of the regular variety $\mathbb P^r$. The height theorem for an ideal
generated by $b-a+1$ elements therefore gives $r\leq b-a+1$,
which is \eqref{a:determinantal-bound}.

In dimension at least three the point kernel is a finite $p$-group
by the already proved A2 reduction. Any nonzero such group contains
an element of order $p$. Applying the preceding exclusion to every
nonzero $p$-torsion line bundle proves the last assertion.
\end{proof}

This supplies a finite linear-algebra test for a fixed pair $(X,L)$
once its $p$-torsion classes and multiplication maps are computed.
The entries depend linearly on the coefficients of the hyperplane.
The full special-hyperplane problem is not settled: the inequality
can hold, and we have not proved that the resulting determinantal
points must lie in the discriminant of singular divisors.
